\documentclass{article}
\usepackage[margin=2cm]{geometry}
\usepackage[dvipsnames]{xcolor}
\usepackage{fontawesome6}
\usepackage{tikz}
\usepackage{tcolorbox}
\usepackage[french]{babel}

\usepackage{fancyhdr} % Added for custom headers

\usetikzlibrary{arrows.meta, positioning, shapes.multipart}

% --- Header Configuration ---
\pagestyle{fancy}
\fancyhf{} % Clears the default header and footer
\lhead{Algorithmes et Pensée Computationelle}
\rhead{École des Sciences Criminelles (UNIL)}
\renewcommand{\headrulewidth}{0.9pt} % Adds a line under the header
% --------------------------


\begin{document}

\noindent
\begin{minipage}[t]{0.39\textwidth}
\vspace{0pt}
\section*{Explication (\texttt{switch} Java \faJava)}
Dans les instructions switch en Java :
\begin{enumerate}
\item Les cas avant la première correspondance sont \textbf{ignorés}
 \item L'exécution commence au \textbf{premier cas correspondant}
 \item Tous les cas suivants sont \textbf{exécutés} indépendamment de leurs étiquettes (fall-through)
 \item L'exécution se poursuit jusqu'à ce qu'une instruction \texttt{break} soit rencontrée
 \item L'instruction \texttt{break} provoque une \textbf{sortie} immédiate du switch
 \item Les cas après le break ne sont \textbf{pas exécutés}
\end{enumerate}
Ce comportement est souvent appelé « fall-through » et est une source courante de bogues lorsque les instructions \texttt{break} sont accidentellement omises.
\end{minipage}
\hspace{0.9cm}
\begin{minipage}[t]{0.3\textwidth}
\vspace{0pt}
\begin{tcolorbox}[width=4.25in,colframe=black, boxrule=1pt, left=0pt, right=0pt, top=0pt, bottom=0pt]
\begin{tikzpicture}[
    node distance=1.5cm,
    % Base style for 2-part nodes (first line + body)
    base/.style={rectangle split, rectangle split parts=2, draw, minimum width=3.2cm, align=center},
    case/.style={base, rectangle split part fill={blue!20,white}},
    matching/.style={base, rectangle split part fill={green!30,white}},
    executed/.style={base, rectangle split part fill={yellow!30,white}},
    break/.style={base, rectangle split part fill={red!30,white}},
    arrow/.style={-Stealth, thick},
    skip/.style={dashed, gray}
]

% Switch cases (first part is header, second part is body)
\node[case] (case1) {\texttt{case expr1}:\nodepart{second}Pas de match};
\node[case, below=of case1] (case2) {\texttt{case expr2}:\nodepart{second}Pas de match};
\node[matching, below=of case2] (case3) {\texttt{case expr3}:\nodepart{second}Cas exécuté\\\textbf{Premier match}};
\node[executed, below=of case3] (case4) {\texttt{case expr4}:\nodepart{second}Cas exécuté\\(Pas de break)};
\node[executed, below=of case4] (case5) {\texttt{case expr5}:\nodepart{second}Cas exécuté\\(Pas de break)};
\node[break, below=of case5] (case6) {\texttt{case expr6}:\nodepart{second}Cas exécuté\\ \texttt{break;}};
\node[case, below=of case6] (case7) {\texttt{case expr\_n (ou default)}:\nodepart{second}Pas exécuté};

% Flow arrows
\draw[skip] (case1) -- (case2);
\draw[skip] (case2) -- (case3);
\draw[skip] (case6) -- (case7);
\draw[arrow, green!60!black, very thick] (case3) -- (case4);
\draw[arrow, orange, very thick] (case4) -- (case5);
\draw[arrow, orange, very thick] (case5) -- (case6);

% Exit arrow
\draw[arrow, red, very thick] (case6.east) -- ++(1.5,0) node[right] {EXIT (à cause du \texttt{break})};

Annotations
\node[right=0.1cm of case1, yshift=1cm] (start) {\texttt{switch (variable ou expression)}};
\draw[arrow] (start) -- (case1);

\node[right=1.5cm of case3, yshift=0.12cm] (match) {\textcolor{OliveGreen}{Exécution commence}};
\draw[arrow, green!60!black] (match) -- ([yshift=0.13cm]case3.east);

\node[right=1.5cm of case4, align=center] (fallthroughA) {Fall-through (tombe-à-travers)\\L'exécution continue};
\draw[arrow, orange] (fallthroughA) -- (case4.east);

\node[right=1.5cm of case5, align=center] (fallthroughB) {Fall-through (tombe-à-travers)\\L'exécution continue};
\draw[arrow, orange] (fallthroughB) -- (case5.east);

\end{tikzpicture}
\end{tcolorbox}
\end{minipage}

\end{document}
